\documentclass{article}
\usepackage{amsmath,amssymb,hyperref}
\title{Cell-functorial retained-state closure on generated native towers}
\date{15 November 2026}
\begin{document}\maketitle
\paragraph{Status and meaning.} Construction for independent review, conditional on the reviewed higher dialgebra and the generated-tower normalization. Closure retains its state as presentation data. It is not an instantaneous trace, does not identify state comparisons with equalities, and does not choose a solution of a feedback equation.

\paragraph{A fixed-state inclusion lemma.} Let $S$ be a bicategory, $K$ a tricategory, and $P,Q:S\to K$ trihomomorphisms. Write $D(P,Q)$ for the native dialgebra bicategory: an object is $(X,f:PX\to QX)$, an arrow is $(h,\alpha:Qh\,f\Rightarrow g\,Ph)$, and a two-cell retains both a state two-cell and the actual process three-cell witnessing its square. For each fixed $X$, there is a pseudofunctor
\[
 C_X:K(PX,QX)\longrightarrow D(P,Q).
\]
It sends $f$ to $(X,f)$. Given a process two-cell $\alpha:f\Rightarrow g$, its native arrow has state $1_X$ and body square
\[
 Q(1_X)f\ \simeq\ f\ \xRightarrow{\alpha}\ g\ \simeq\ gP(1_X).
\]
The first structural comparison uses the inverse of the chosen $Q$ unit equivalence, the second the forward $P$ unit equivalence, with the actual tricategory unitors. A process three-cell $\Gamma:\alpha\Rightarrow\alpha'$ maps to the identity two-cell on $1_X$ and the three-cell obtained by whiskering $\Gamma$ with exactly these unit comparisons. The pseudofunctorial action of $P,Q$ on identity state two-cells supplies the remaining identity frames. This does not erase or invert $\Gamma$.

\paragraph{Proof of the local laws.} Composing two image arrows first produces state $1_X1_X$, not literally $1_X$. Its comparison with the image of the composite uses the state unitor and the $P,Q$ multiplication/unit comparison modifications. Expand both native body squares: their chronological nonstructural word is $\alpha$ followed by $\beta$. The adjacent unit frames cancel through the chosen adjunction triangles, leaving the framed composite. This defines an invertible native two-cell. The identity comparison is the analogous empty-word paste. Naturality at actual $\Gamma,\Delta$ follows by whiskering and tricategorical interchange. For three arrows both associativity pastes reduce to the same chronological word with the specified $P,Q$ unit/coherence modifications; their equality is the trihomomorphism unit coherence and the tricategory triangle equations. The identity laws use the same modifications. No equality of state components is used as a substitute for equality of the process three-cell witnesses.

\paragraph{Closure at every generated finite depth.} Put $K=T_n$, the independently constructed normalized level from the original process bicategory $K_0$, and take $P_A(X)=JX\otimes A$, $Q_B(X)=JX\otimes B$ with their unit-state trihomomorphisms. Follow $C_X$ by the raw-to-normalized hom biequivalence. This constructs
\[
 C^n_{X;A,B}:T_n(JX\otimes A,JX\otimes B)
 \longrightarrow T_{n+1}(A,B)
\]
as a pseudofunctor, not merely a function on machines. A machine with inner coordinates $(X_2,\ldots,X_{n+1})$ acquires the new outer coordinate $X_1=X$. The body is conjugated by the specified action equivalence identifying its domain with
$X_{n+1}\otimes\cdots\otimes X_1\otimes A$, and similarly at the output. A native process arrow acquires the outer identity state map; all inner state maps and its body square remain. A native process three-cell acquires the outer identity two-cell; its inner tuple and its actual equation remain. This construction uses only the individual $P_A,Q_B$ and the established hom normalization, not a complete action doctrine on arbitrary process tricategories.

\paragraph{Nested closure and finite joining.} For a body with boundary $JY\otimes(JX\otimes A)$ and the analogous output, closing $Y$ and then $X$ retains the ordered coordinates $(X,Y,\mathbf Z)$; its total state is $m_n(\mathbf Z)\otimes Y\otimes X$. It is not a same-state fusion. The already specified joining map replaces the consecutive pair $(X,Y)$ by $Y\otimes X$, with its action associativity comparisons. On process arrows and three-cells this comparison is obtained by the same unit-frame insertion, followed by the strong monoidal state-base-change comparison. The required finite joining coherence is therefore the existing block-joining coherence, provided the generated normalization comparison is used throughout. This statement does not create an inverse effectful interchange or a new quotient by sliding.

\paragraph{What is deliberately not concluded.} An arbitrary state map $h:X\to Y$ alone does not produce a comparison between closing unrelated bodies. Such a native comparison additionally requires a square $Qh\,f\Rightarrow g\,Ph$ and, for higher comparisons, its actual three-cell witness. Nor does the existence of these local closure pseudofunctors establish a universal evaluation algebra, full Conway identities, or compatibility with arbitrary shared-interface synchronization. Those require their separately typed structural laws. The present result makes explicit that the generated finite tower supports closure functorially on all its admitted native cells, rather than only on objects or decategorified behavior.
\end{document}
